\chapter*{Словарь терминов}             % Заголовок
\addcontentsline{toc}{chapter}{Словарь терминов}  % Добавляем его в оглавление

\textbf{ab initio} : квантово-химический расчет из первых принципов, выполняемый без подгонки параметров непосредственно под анализируемый экспериментальный спектр.

\textbf{Bruker IFS~125HR} : Фурье-спектрометр высокого разрешения, использованный для регистрации экспериментальных инфракрасных спектров.

\textbf{\(C_{2v}\)} : точечная группа симметрии молекулы сероводорода H$_2$S и симметричных изотопологов типа D$_2$S.

\textbf{\(C_s\)} : точечная группа симметрии, содержащая одну плоскость симметрии; в работе используется при описании молекулы HDS.

\textbf{\(d_{\mathrm{rms}}\)} : среднеквадратичное отклонение экспериментальных значений от рассчитанных.

\textbf{\(\mathrm{D_2S}\)} : дидейтерированный сероводород, изотополог H$_2$S, в котором оба атома водорода заменены атомами дейтерия.

\textbf{GF} : метод Вильсона для расчета нормальных колебаний молекул; \(G\) обозначает матрицу кинематических коэффициентов, \(F\) --- силовую матрицу.

\textbf{\(\mathrm{GeH_4}\)} : герман, тетрагидрид германия.

\textbf{HDS} : монодейтерированный сероводород; в работе также рассматриваются изотопические разновидности HD$^{32}$S и HD$^{34}$S.

\textbf{ИК} : инфракрасная область спектра; в тексте используется в выражениях <<ИК-спектры>> и <<ИК-диапазон>>.

\textbf{\(\mathrm{KBr}\)} : бромид калия, материал оптических элементов, указанный в условиях регистрации спектров.

\textbf{MCT} : детектор на основе теллурида ртути--кадмия (mercury cadmium telluride), применяемый при регистрации инфракрасных спектров.

\textbf{\(\mathrm{SiH_4}\)} : силан, тетрагидрид кремния.

\textbf{STDS} : Spherical Top Data System, база данных и программный комплекс для расчета и анализа спектров молекул типа сферического волчка.

\textbf{\(T_d\)} : точечная группа тетраэдрической симметрии, используемая при описании молекул типа XY$_4$.

\textbf{XY$_2$} : обобщенная формула трехатомной молекулы с центральным атомом X и двумя атомами Y.

\textbf{XY$_4$} : обобщенная формула тетраэдрической молекулы с центральным атомом X и четырьмя атомами Y.
